% ===========================================================================
\chapter{Record and replay}\label{ch:replay}

A real model is a nuisance to test with. It costs money, it needs the
network, and it never answers the same way twice. That last part matters
most: you can't write a check that says ``after this request, the reply is
that'', and you can't show a session in a tutorial and expect a reader to
get the same one.

The usual answer, in any program that talks to a service, is to
\emph{record} a real session once, and \emph{replay} it afterwards. While
recording, every request goes to the real model, and every reply is saved
to a file along with what was sent. While replaying, the replies come from
the file instead. The agent can't tell the difference: it is given a
\code{Model}, and both are models.

A replay is only faithful if the agent does the same things it did when
the session was recorded. The replies are fixed; but the agent really runs
the tools, so the results it sends back are produced afresh each time, from
the workspace as it is. If the sample project were different, or
\code{nex} printed its errors differently, the results would differ, and
the recorded replies (which answered the old results) might no longer make
sense. So the replay compares what is sent with what was sent when the
session was recorded, and notes where they first differ.

\section{The session file}

A session is a JSON file holding the task, the folder the workspace was
copied from, the model's name, and the turns. Each turn is the new part of
one request (the last user message, from Chapter~\ref{ch:model}) and the
model's response, exactly as it came:

\begin{json}
{"model": "claude-opus-5-5",
 "task": "Some of the checks fail. Find out why and fix it.",
 "template": "sample_project",
 "turns": [{"sent": [{"type": "text", "text": "Some of the checks fail. ..."}],
            "response": {"content": [...], "stop_reason": "tool_use", ...}},
           {"sent": [{"type": "tool_result", ...}],
            "response": ...},
           ...]}
\end{json}

Both new models read and write files, so \code{model.nex} needs
\code{io/Path}:

\begin{tightaddbefore}{model.nex}{after intern Messages}{|-- Everything the model is given for one reply.}
intern io/Path
\end{tightaddbefore}

\section{Playing back}

\code{Replay_Model} reads a session, and gives its responses back in order.
Before each one, it compares the request it was sent with the recording.
\code{diverged_at} is the first reply whose request was different, or 0 if
there has been none.

\begin{append}{model.nex}
-- A model that plays back a session recorded by Recording_Model: the
-- same replies, in the same order, with no network and no cost.
--
-- A recorded reply only makes sense after what was sent before it, so
-- it compares each request with the recording. If the tools gave
-- different results this time (the project's files are not the ones
-- the session was recorded with, say), the replies may no longer fit,
-- and it says so.
class Replay_Model inherit Model
create
  -- The session saved in the file `path`.
  make(path: String)
  require
    file_exists: (create Path.make(path)).is_file()
  do
    session := json_value((create Path.make(path)).read_text())
    diverged_at := 0
  end

feature
  session: Json_Value
  played: Integer        -- how many replies it has given
  diverged_at: Integer   -- the first reply whose request differed, or 0

  name(): String do
    result := session.get("model").as_string + " (recorded)"
  end

  -- The task the session was recorded with.
  task(): String do
    result := session.get("task").as_string
  end

  -- The folder the session's workspace was copied from.
  template(): String do
    result := session.get("template").as_string
  end

  turns(): Integer do
    result := session.get("turns").length
  end

  can_reply(): Boolean do
    result := played < turns()
  end

  reply(request: Request): Reply do
    let turn: Json_Value := session.get("turns").at(played)
    played := played + 1
    if diverged_at = 0 and last_message_json(request) /= turn.get("sent").to_string then
      diverged_at := played
    end
    result := create Reply.from_json(turn.get("response"))
  end
end
\end{append}

A replay that diverges carries on regardless. Often the difference is
harmless (a line in an error message that changed between versions of
\code{nex}), and the session still plays out sensibly; and if it doesn't,
seeing how it goes wrong is more useful than a refusal to play. The program
will say so at the end.

\section{Recording}

\code{Recording_Model} is a \emph{wrapper}: it is given another model, and
passes every request on to it. Each time a reply comes back, it adds the
turn to the session, and saves the file. Saving after every reply, rather
than at the end, means that a session stopped part way (with Ctrl-C,
say) is still kept, as far as it went.

\begin{append}{model.nex}
-- A model that passes every request on to another model and saves each
-- exchange to a file as it goes, for Replay_Model to play back.
class Recording_Model inherit Model
create
  make(inner: Model, path, task, template: String) do
    this.inner := inner
    this.path := path
    this.task := task
    this.template := template
  end

feature
  inner: Model
  path: String         -- where the session is saved
  task: String
  template: String
  turns: Array[Any]

  name(): String do
    result := inner.name()
  end

  can_reply(): Boolean do
    result := inner.can_reply()
  end

  reply(request: Request): Reply do
    result := inner.reply(request)
    let turn: Map[String, Any] := {
      "sent": json_value(last_message_json(request)).raw(),
      "response": result.response.raw()
    }
    turns.add(turn)
    save()
  end

  private feature
    -- Saved after every reply, so a session stopped part way is kept.
    save() do
      let session: Map[String, Any] := {"model": inner.name(), "task": task,
        "template": template, "turns": turns}
      (create Path.make(path)).write_text(json_text(session))
    end
end
\end{append}

\code{Recording_Model} can reply whenever the model inside it can: that is
the \code{can_reply} it redefines. (Without it, a recording of a script
would claim it could always reply, and break \code{reply}'s precondition
when the script ran out.)

This is a pattern worth knowing. A wrapper adds something (here,
recording) to any model, without changing it and without the agent
knowing. The same trick would log every request, or count the cost, or
stop after a budget.

\section{Trying it}

We can record a scripted session, play it back, and then play it back
again on a workspace that has changed:

\begin{tryit}{try_replay.nex}
intern Agent

let task: String := "What does fines.nex do?"
let quiet: Reporter := create Reporter.make(true, false)
let ws: Workspace := fresh_workspace(user_path("sample_project"), user_path("scratch"))
let script: Scripted_Model := create Scripted_Model.make([
  use_tool("", "read_file", '{"path": "fines.nex"}'),
  say("It charges a fine per day late, up to a maximum.")
])
let recorder: Recording_Model := create Recording_Model.make(script,
  user_path("session.json"), task, "sample_project")
create Agent.make(recorder, standard_tools(ws), create Allow_Everything.make(), quiet).run(task)

let replay: Replay_Model := create Replay_Model.make(user_path("session.json"))
console.print_line("the session has " + count_text(replay.turns(), "turn")
  + "; its task: " + replay.task())
let again: Run_Result := create Agent.make(replay, standard_tools(ws),
  create Allow_Everything.make(), quiet).run(replay.task())
console.print_line(again.ending + ": " + again.text + " (diverged at "
  + replay.diverged_at.to_string + ")")

ws.write("fines.nex", "-- no fines any more\n")
let changed: Replay_Model := create Replay_Model.make(user_path("session.json"))
create Agent.make(changed, standard_tools(ws), create Allow_Everything.make(), quiet).run(task)
console.print_line("after the change, diverged at " + changed.diverged_at.to_string)
\end{tryit}
\tryoutput{try_replay}

The first replay matched the recording all the way. The second was sent a
different \code{fines.nex}, so its second request (which carried the file)
differed from the recording, and the replay said so.

\begin{aside}
\textbf{The session that comes with this tutorial.}
\code{sessions/fix_late_fines.json} is the session played back in
Chapter~1 and in the next chapter. Its \code{model} field says who played
the model's part, and the program prints it when the session is played.
With \code{--record} (next chapter), the agent records sessions of your
own.
\end{aside}
